% =====================================================================
%  Presentaciones con Beamer (a fondo) — suite de documentos LaTeX
%  Nivel 3 (Tipos de documento) · clase article (guía SOBRE la clase beamer).
%  Compilar: latexmk -pdf beamer.tex
%  (No es la Masterclass: este doc trabaja la clase beamer en profundidad.)
% =====================================================================
\documentclass[11pt,a4paper]{article}

% --- Base estándar ---
\usepackage[utf8]{inputenc}
\usepackage[T1]{fontenc}
\usepackage{lmodern}
\usepackage{microtype}
\usepackage[spanish,es-noshorthands]{babel}
\usepackage{csquotes}
\usepackage{amsmath}\usepackage{amssymb}
\usepackage{booktabs}
\usepackage{tabularx}
\usepackage[margin=2.4cm]{geometry}
\usepackage{enumitem}
\usepackage{xcolor}
\usepackage{listings}
\usepackage{hyperref}
\hypersetup{unicode, pdftitle={Presentaciones con Beamer (a fondo)},
            pdfauthor={Rodrigo Martín Vidal Galán},
            colorlinks=true, linkcolor=blue!50!black, urlcolor=blue!50!black}
\usepackage[spanish,capitalise,nameinlink]{cleveref}

% --- Estilo de código ---
\definecolor{codeback}{RGB}{245,247,250}
\definecolor{codekw}{RGB}{20,60,120}
\definecolor{codecom}{RGB}{110,120,135}
\lstset{
  basicstyle=\ttfamily\footnotesize,
  keywordstyle=\color{codekw}\bfseries,
  commentstyle=\color{codecom}\itshape,
  backgroundcolor=\color{codeback},
  frame=single, rulecolor=\color{codekw!30},
  breaklines=true, showstringspaces=false, columns=fullflexible,
  language=[LaTeX]TeX, xleftmargin=8pt, framesep=4pt, extendedchars=true,
  literate={á}{{\'a}}1 {é}{{\'e}}1 {í}{{\'i}}1 {ó}{{\'o}}1 {ú}{{\'u}}1
           {ñ}{{\~n}}1 {¿}{{?`}}1 {¡}{{!`}}1,
}
\newcommand{\C}[1]{\texttt{\textbackslash #1}}
\newcommand{\pkg}[1]{\textsf{#1}}

\title{\textbf{Presentaciones con Beamer (a fondo)}\\[2pt]
       \large La clase \texttt{beamer} en profundidad}
\author{Rodrigo Martín Vidal Galán}
\date{Junio de 2026}

\begin{document}
\maketitle

\noindent Guía de referencia de la clase \pkg{beamer}, la forma estándar de hacer
\textbf{presentaciones} con \LaTeX. Acá se explica \emph{cómo funciona} cada
pieza —frames, overlays, bloques, temas— para que puedas armar y personalizar tus
diapositivas. \emph{(Es distinto de la \emph{Masterclass}, que es una charla
\textbf{sobre LaTeX en general}; este documento entra a fondo en la clase
\texttt{beamer}.)}

\newpage
\tableofcontents
\newpage

% =====================================================================
\section{Anatomía de una presentación}
% =====================================================================
Una presentación es una serie de \textbf{frames} (diapositivas). El esqueleto
mínimo:
\begin{lstlisting}
\documentclass[10pt,aspectratio=169]{beamer}
\usetheme{Madrid}
\begin{document}
  \begin{frame}{Titulo del frame}
    Contenido...
  \end{frame}
\end{document}
\end{lstlisting}
\begin{itemize}[noitemsep,leftmargin=*]
  \item \texttt{aspectratio=169} da pantalla ancha 16:9 (también \texttt{43},
        \texttt{1610}).
  \item Todo lo \textbf{visible} vive dentro de un \texttt{frame}.
  \item \pkg{beamer} ya carga \pkg{hyperref}, \pkg{xcolor}, \pkg{amsmath} y
        \pkg{tikz}: \textbf{no} los recargues.
\end{itemize}

\paragraph{El frame y sus opciones.} El entorno \texttt{frame} acepta un título
(y subtítulo) y opciones entre corchetes:
\begin{lstlisting}
\begin{frame}[fragile]{Titulo}{Subtitulo opcional}
  ...
\end{frame}
\end{lstlisting}
\begin{center}\small
\renewcommand{\arraystretch}{1.3}
\begin{tabularx}{\linewidth}{@{}lX@{}}
\toprule
\textbf{Opción} & \textbf{Para qué} \\
\midrule
\texttt{fragile} & \textbf{Obligatoria} si el frame contiene \texttt{verbatim}/
                   \pkg{listings} (código). \\
\texttt{plain}   & Sin encabezado ni pie (portada, o una imagen a sangre). \\
\texttt{allowframebreaks} & Si el contenido no entra, lo parte en varias diapositivas. \\
\texttt{t} / \texttt{c} / \texttt{b} & Alinea el contenido arriba / centro / abajo. \\
\bottomrule
\end{tabularx}
\end{center}

% =====================================================================
\section{Estructura y navegación}
% =====================================================================
\C{section} y \C{subsection}, puestos \textbf{fuera} de los frames, organizan la
charla y alimentan el índice:
\begin{lstlisting}
\section{Introduccion}
\begin{frame}{...} ... \end{frame}

% un indice que se actualiza solo:
\begin{frame}{Contenido}\tableofcontents\end{frame}
\end{lstlisting}
Para mostrar un \textbf{mini-índice} al empezar cada sección, resaltando la
actual, se automatiza con \C{AtBeginSection}:
\begin{lstlisting}
\AtBeginSection[]{
  \begin{frame}{Contenido}
    \tableofcontents[currentsection]
  \end{frame}}
\end{lstlisting}
Así el público siempre sabe \enquote{dónde está} en la presentación.

% =====================================================================
\section{Bloques}
% =====================================================================
Los \textbf{bloques} encuadran contenido con un título y un color, según su
\emph{significado}. Hay tres:
\begin{lstlisting}
\begin{block}{Definicion}        Texto neutro.        \end{block}
\begin{alertblock}{Cuidado}      Algo importante.     \end{alertblock}
\begin{exampleblock}{Ejemplo}    Un caso concreto.    \end{exampleblock}
\end{lstlisting}
El \texttt{block} usa el color de estructura del tema (azul, en Madrid);
\texttt{alertblock} sale en rojo (para advertencias) y \texttt{exampleblock} en
verde (para ejemplos). La idea es elegirlos por lo que \textbf{significan}, no por
el color que producen.

% =====================================================================
\section{Columnas}
% =====================================================================
Para poner cosas lado a lado (típico: texto a un lado, gráfico al otro) está el
entorno \texttt{columns}:
\begin{lstlisting}
\begin{columns}[T]   % [T]=alinear arriba; [c]=centrar verticalmente
  \begin{column}{0.48\textwidth}
    Texto, items o ecuaciones...
  \end{column}
  \begin{column}{0.48\textwidth}
    \includegraphics[width=\textwidth]{figura}  % o un tikzpicture
  \end{column}
\end{columns}
\end{lstlisting}
\begin{quote}\small
\textbf{Regla práctica:} la suma de los anchos \emph{más} la separación debe ser
$\le \texttt{\textbackslash textwidth}$. Si te pasás (p.\,ej.\ dos columnas de
\texttt{0.5} con separación), las columnas se van al margen.
\end{quote}

% =====================================================================
\section{Overlays: revelado por pasos}
% =====================================================================
Lo que distingue a \pkg{beamer}: un mismo frame puede mostrarse en \textbf{varios
pasos} (cada paso es una \enquote{diapositiva} del PDF). La forma más simple es
\C{pause}:
\begin{lstlisting}
Primero esto. \pause Despues esto.
\end{lstlisting}
Para más control se usan \textbf{especificaciones de overlay} \texttt{<...>}, que
indican en qué pasos aparece cada cosa:
\begin{lstlisting}
\begin{itemize}
  \item<1-> Visible desde el paso 1.
  \item<2-> Aparece en el paso 2.
  \item<3-> Y este en el 3.
\end{itemize}
\end{lstlisting}

\paragraph{La familia de comandos de overlay.}
\begin{center}\small
\renewcommand{\arraystretch}{1.3}
\begin{tabularx}{\linewidth}{@{}lX@{}}
\toprule
\textbf{Comando} & \textbf{Efecto} \\
\midrule
\C{only<n>\{\dots\}}   & Muestra el material \textbf{solo} en el paso $n$; \emph{no} reserva espacio. \\
\C{uncover<n->\{\dots\}} & Lo revela en $n$, pero \textbf{reserva el lugar} desde el inicio. \\
\C{onslide<n->\{\dots\}} & Como \texttt{uncover}, versión flexible. \\
\C{alt<n>\{A\}\{B\}}   & Muestra \texttt{A} en el paso $n$ y \texttt{B} en el resto. \\
\C{temporal<n>\{antes\}\{durante\}\{después\}} & Tres estados según el paso. \\
\C{alert<n>\{texto\}}  & \textbf{Resalta} el texto solo en el paso $n$. \\
\bottomrule
\end{tabularx}
\end{center}
\textbf{Sutileza clave:} \C{only} \emph{no} reserva espacio (lo de abajo se mueve
al aparecer/desaparecer); \C{uncover} sí. Para que el contenido no \enquote{salte}
entre pasos, preferí \C{uncover}.

Los overlays también funcionan \textbf{dentro} de ecuaciones y de gráficos
\pkg{pgfplots} (p.\,ej.\ \C{onslide<2->\{\C{addplot}\dots\}} dibuja una curva
recién en el segundo paso), lo que permite \enquote{construir} una figura por
etapas.

% =====================================================================
\section{Temas y personalización}
% =====================================================================
El aspecto se arma combinando \textbf{cuatro temas}:
\begin{center}\small
\renewcommand{\arraystretch}{1.3}
\begin{tabularx}{\linewidth}{@{}lX@{}}
\toprule
\textbf{Tema} & \textbf{Controla} \\
\midrule
\texttt{theme}      & El \enquote{maestro} (\texttt{Madrid}, \texttt{Berlin}, \texttt{metropolis}\dots). \\
\texttt{colortheme} & La paleta (\texttt{whale}, \texttt{beaver}, \texttt{seagull}\dots). \\
\texttt{innertheme} & Viñetas, bloques (\texttt{circles}, \texttt{rectangles}). \\
\texttt{outertheme} & Cabecera, pie, barra lateral. \\
\bottomrule
\end{tabularx}
\end{center}
\begin{lstlisting}
\usetheme{Madrid}\usecolortheme{whale}
\useinnertheme{circles}
\end{lstlisting}
\begin{quote}\small
\textbf{Nota (2026):} \texttt{metropolis} (tema moderno y minimalista) está
\textbf{sin mantenimiento}; su fork \texttt{moloch} lo continúa.
\end{quote}

\paragraph{Retocar piezas sueltas.} Sin cambiar de tema, podés ajustar colores,
fuentes y plantillas concretas:
\begin{lstlisting}
\setbeamercolor{frametitle}{bg=blue!70!black, fg=white}
\setbeamerfont{frametitle}{series=\bfseries}
\setbeamertemplate{navigation symbols}{}   % oculta la botonera de navegacion
\setbeamertemplate{caption}[numbered]       % numera las figuras
\setbeamertemplate{itemize item}[triangle]
\end{lstlisting}
\begin{itemize}[noitemsep,leftmargin=*]
  \item \C{setbeamercolor} $\to$ colores de un elemento (\texttt{fg}/\texttt{bg}).
  \item \C{setbeamerfont} $\to$ tamaño / serie / forma.
  \item \C{setbeamertemplate} $\to$ \textbf{reemplaza} una plantilla entera.
\end{itemize}

% =====================================================================
\section{Matemática y gráficos en las diapositivas}
% =====================================================================
Toda la matemática de \pkg{amsmath} funciona igual que en un \texttt{article}.
Lo propio de \pkg{beamer} es combinarla con overlays para \textbf{construir} una
idea por pasos:
\begin{lstlisting}
\begin{frame}{Tangente}
  \begin{columns}[c]
    \begin{column}{0.5\textwidth}
      \begin{itemize}
        \item<1-> La parabola $y=x^2$...
        \item<2-> ...y su recta tangente en $x=1$.
      \end{itemize}
    \end{column}
    \begin{column}{0.5\textwidth}
      \begin{tikzpicture}
        \begin{axis}[axis lines=middle, domain=-0.5:2, samples=60]
          \addplot[blue,thick]{x^2};
          \onslide<2->{\addplot[red,thick,domain=0:2]{2*x-1};}
        \end{axis}
      \end{tikzpicture}
    \end{column}
  \end{columns}
\end{frame}
\end{lstlisting}
En el primer paso se ve la parábola; en el segundo aparece la tangente (porque su
\C{addplot} está envuelto en \C{onslide<2->}).

% =====================================================================
\section{Notas del orador y modo entrega}
% =====================================================================
\paragraph{Notas.} \C{note\{\dots\}} agrega notas que solo ve quien presenta (en
el monitor de presentador):
\begin{lstlisting}
\note{Recordar: mencionar el contraejemplo.}
\setbeameroption{show notes on second screen=right}
\end{lstlisting}

\paragraph{Imprimir / repartir.} Para una versión sin animaciones (cada frame
colapsado a una sola diapositiva) y con varias por hoja:
\begin{lstlisting}
\documentclass[handout]{beamer}   % colapsa todos los overlays
% varias diapositivas por hoja (con pgfpages):
\usepackage{pgfpages}
\pgfpagesuselayout{4 on 1}[a4paper,landscape,border shrink=5mm]
\end{lstlisting}
La opción \texttt{handout} es clave: convierte una presentación llena de
\C{pause} en un PDF limpio para entregar.

% =====================================================================
\section{Buenas prácticas}
% =====================================================================
\begin{itemize}[noitemsep,leftmargin=*]
  \item \textbf{Una idea por diapositiva}; poco texto, frases cortas.
  \item Overlays \textbf{con moderación}: revelar para guiar la atención, no por
        efecto.
  \item Preferí \C{uncover} a \C{only} para que el contenido no \enquote{salte}.
  \item \texttt{[fragile]} en \textbf{todo} frame con código.
  \item Bloques por su \emph{significado} (\texttt{block}/\texttt{alertblock}/
        \texttt{exampleblock}), no por su color.
  \item Generá un \texttt{handout} para imprimir o repartir.
  \item Anchos de columna $+$ separación $\le \texttt{\textbackslash textwidth}$.
\end{itemize}

% =====================================================================
\section*{Ver también}
% =====================================================================
\addcontentsline{toc}{section}{Ver también}
\begin{itemize}[noitemsep,leftmargin=*]
  \item \emph{Masterclass de LaTeX} — panorama general (esta es la versión
        \emph{a fondo} de la clase \texttt{beamer}).
  \item \emph{Plantillas} — \texttt{beamer.tex} para arrancar una presentación.
  \item \emph{Gráficos} — \pkg{pgfplots} dentro de los frames.
  \item Manual: \texttt{texdoc beamer} (la guía oficial, exhaustiva).
\end{itemize}

\end{document}
